\documentclass[11pt]{article}

\usepackage[margin=25mm]{geometry}

\usepackage{amsmath,amssymb,fontspec,unicode-math}

\IfFontExistsTF{DejaVu Serif}{\setmainfont{DejaVu Serif}}{\setmainfont{Latin Modern Roman}}

\IfFontExistsTF{Latin Modern Math}{\setmathfont{Latin Modern Math}}{\setmathfont{DejaVu Math TeX Gyre}}

\usepackage{xurl,hyperref}

\setlength{\parindent}{0pt}\setlength{\parskip}{6pt}

\begin{document}

\title{GeRaF: Neural Geometry Reconstruction from Radio Frequency Signals}

\author{ELVIA preliminary AI analysis}\date{21 October 2026}\maketitle

Preliminary assessment: Major.\par

Independent scientific review remains pending. Author adoption has not been requested.

The central additive lensless-transmittance derivation fails even on a valid monotone planar SDF. The printed matched-filter derivative also lacks a resultant factor, but the later release supplies a correct magnitude-chain path.

\section*{Review scope}

Latest-version, claim-focused review of the cited author-paper passages and applicable current author-release paths. The stated counterexamples and scope controls were checked analytically; no author program, old verifier, training run or benchmark was executed.

\section*{Problem 1: Additive calibration conflicts with the defining transmission integral}

Source: Current arXiv v2 Section 6, PDF p.7, Eq.(7); Appendix B.1–B.2, PDF pp.21–22

Assessment: supported; impact: major

\subsection*{Source claim (paraphrase)}

Section 6, Eq.(7), uses \(T_r=T_p-C_p(0)+C_r(0)\) to calibrate reused-ray transmittance. Appendix B defines density from the clipped logarithmic derivative of the logistic CDF \(C\).

\subsection*{Analysis}

On a monotone entering ray the integral gives \(T(u)=C(u)/C(0)\). Use the true planar SDF \(f=(x+z)/\sqrt2\) and unit logistic scale. Origins \((0,0,\sqrt2\log3)\) and \((-\sqrt2\log2,0,\sqrt2\log3)\) lie in the same aperture plane and have CDFs \(3/4\), \(3/5\); choose their common endpoint at CDF \(1/2\). Then \(T_p=2/3\), but additive calibration gives \(31/60\), while the receiver ray requires \(5/6\). All values are finite and positive, and both straight rays are monotone for this planar SDF.

\subsection*{Finding and scope}

The central calibration formula is not derived from the stated density, even before nonmonotone path complications. No reconstruction benchmark is remeasured.

\subsection*{Proposed correction}

For monotone entering rays, the correct origin adjustment is \(T_r=T_pC_p(0)/C_r(0)\). General paths instead depend on accumulated downward variation of \(\log C\), so this endpoint ratio is not a universal ray-reuse repair.

An explicit local correction is described; author adoption and a corresponding published erratum are not established.

\section*{Problem 2: The printed magnitude derivative omits the complex resultant}

Source: Current arXiv v2 Appendix A.5, PDF p.19, Eq.(13); Algorithm 1, PDF p.20

Assessment: supported; impact: minor

\subsection*{Source claim (paraphrase)}

Appendix A.4 differentiates \(P_j=|z_j|\) but omits the resultant needed in its magnitude derivative.

\subsection*{Analysis}

For one positive real input with coefficient \(1\) and signal \(2\), \(P=2\), true derivative is \(1\), while the printed rule gives \(1/P=1/2\). This real subcase avoids complex-gradient convention ambiguity. The later `rf\_rendering.py` instead computes real/imaginary outputs followed by `torch.hypot`, which retains the needed chain rule.

\subsection*{Finding and scope}

The printed formula is wrong; a current released gradient defect and historical benchmark effect are not established.

\subsection*{Proposed correction}

For \(z\ne0\), use \(d|z|=\operatorname{Re}(\bar z\,dz)/|z|\), with a declared convention for the final real-loss complex gradient and the zero boundary.

An explicit local correction is described; author adoption and a corresponding published erratum are not established.

\section*{Limits}

Preliminary AI assessment; independent human scientific adjudication remains pending.\par

An analytic counterexample or a release discrepancy does not by itself remeasure the paper’s reported experiments.\par

The later GeRaF-SENS code and later arXiv extension are distinct from the original proceedings implementation history.\par

\section*{Source identity}

arXiv:2605.29097v2 (current author version; updated 2026-08-04T20:21:35Z)

\url{https://arxiv.org/pdf/2605.29097v2}

Reviewed PDF SHA-256: \nolinkurl{363ff65c4a39903188162ec734e9e40f523feeba6454880b279ec41d8b01ccaa}

The selected author/publisher source is the same work as the original case, matched by exact title, author identity and official version channel. The exact selected PDF SHA-256 was recomputed in the latest-version pass for Issue \#27. The PDF is a retained author/publisher copy; its original retrieval time is stated only when supported by matching acquisition metadata. File modification time was not substituted.

\end{document}
